\chapter*{Введение}
\addcontentsline{toc}{chapter}{Введение}

Музыкальное сопровождение заведений и мероприятий является важной частью клиентского опыта. В ресторанах, кафе, барах, гостиницах, салонах красоты, фитнес-центрах и event-пространствах музыка не только заполняет звуковой фон, но и участвует в формировании атмосферы, восприятия бренда и эмоционального состояния посетителей. При этом в большинстве практических сценариев музыкальный поток формируется заранее: через статичные плейлисты, ручной выбор администратора, DJ-сет или специализированный B2B-сервис. Такая модель позволяет контролировать общий стиль, но почти не учитывает фактический состав аудитории, находящейся в пространстве в конкретный момент времени.

Актуальность темы связана с несколькими тенденциями. Во-первых, пользователи музыкальных стримингов уже привыкли к персонализированным рекомендациям и ожидают, что цифровые сервисы учитывают их вкусы. Во-вторых, бизнес все чаще рассматривает музыку как элемент управления клиентским опытом, а не как случайный фон. В-третьих, для публичного воспроизведения музыки в России важны правовая модель, взаимодействие с РАО и ВОИС, а также отчетность о фактически использованных произведениях и фонограммах. В-четвертых, развитие рекомендательных систем и наличие крупных музыкальных датасетов, в том числе YAMBDA~\cite{yambda2025}, позволяют экспериментально исследовать методы формирования групповых музыкальных рекомендаций.

Проблема исследования состоит в том, что между персональными музыкальными рекомендациями и практикой публичного музыкального сопровождения существует разрыв. Персональные сервисы хорошо учитывают вкусы отдельного пользователя, но не предназначены для коммерческого публичного воспроизведения. B2B-сервисы музыки для бизнеса помогают легально и централизованно управлять музыкальным фоном, но обычно ориентируются на тип заведения, заранее заданный плейлист или брендовый стиль, а не на текущую аудиторию. Сервисы заказа музыки гостями добавляют интерактивность, однако дают непропорциональное влияние отдельным пользователям и могут нарушать целостность атмосферы. Поэтому представляется перспективным исследовать модель, в которой музыкальный поток формируется как коллективная рекомендация для присутствующей группы посетителей, но остается ограниченным рамками бизнеса.

Объектом исследования является процесс формирования музыкального сопровождения заведений и мероприятий с использованием данных о предпочтениях пользователей музыкального сервиса. Предметом исследования являются продуктовая концепция и алгоритмические подходы к генерации коллективных музыкальных рекомендаций для группы посетителей в условиях ограничений публичного воспроизведения, приватности и бизнес-контекста.

Цель работы~--- разработать и обосновать концепцию сервиса коллективных музыкальных рекомендаций для заведений и мероприятий, а также исследовать техническую возможность формирования групповых музыкальных рекомендаций на основе открытого датасета YAMBDA. В рамках работы сервис рассматривается как потенциальная часть экосистемы «Яндекс Музыки»: для бизнеса создается специальный профиль, пользователи добровольно участвуют в формировании общей музыкальной «волны», а обработка персональных предпочтений остается внутри платформенного контура.

Для достижения цели поставлены следующие задачи:

\begin{enumerate}
    \item проанализировать предметную область музыкального сопровождения заведений и мероприятий;
    \item описать сегменты рынка, текущие практики использования музыки и существующие B2B-решения;
    \item рассмотреть правовые аспекты публичного воспроизведения музыки и требования к отчетности;
    \item выявить потребности и ограничения целевых сегментов;
    \item оценить рыночный потенциал сервиса отдельно для постоянных B2B-точек и событийного рынка;
    \item разработать концепцию сервиса, его ценностное предложение, пользовательские сценарии и модель монетизации;
    \item систематизировать подходы к персональным музыкальным и последовательным рекомендациям;
    \item рассмотреть методы групповых рекомендаций и определить релевантные архитектуры для сравнения;
    \item реализовать и исследовать модель коллективных музыкальных рекомендаций на датасете YAMBDA;
    \item разработать демонстрационный сценарий применения модели и определить направления дальнейшего развития технологии.
\end{enumerate}

Исследовательская значимость работы состоит в соединении продуктовой постановки с технической задачей групповых музыкальных рекомендаций. Работа не ограничивается алгоритмическим экспериментом и рассматривает, в каком рыночном, правовом и пользовательском контексте такая модель может быть применена. Практическая значимость заключается в проработке концепции B2B/B2C-сервиса, который может повысить управляемость музыкальной атмосферы, снизить административную нагрузку бизнеса и предложить пользователям новый приватный способ влияния на музыку в физическом пространстве.

Работа имеет смешанный исследовательско-проектный характер и состоит из шести глав. В главе~\ref{ch:market} анализируются рынок и предметная область музыкального сопровождения заведений и мероприятий, конкурентная среда, правовые ограничения, потребности сегментов и рыночный потенциал сервиса. В главе~\ref{ch:service-concept} разрабатывается концепция сервиса «Резонанс» / «Яндекс Музыка Резонанс», включая целевые аудитории, функции, пользовательские сценарии, позиционирование и монетизацию. В главе~\ref{ch:sequential} рассматриваются теоретические и технологические основы построения музыкальных рекомендательных систем. В главе~\ref{ch:groups} систематизируются подходы к формированию коллективных рекомендаций. В главе~\ref{ch:experiment} описывается реализация и экспериментальное исследование модели коллективных музыкальных рекомендаций на датасете YAMBDA. В главе~\ref{ch:prototype} рассматривается демонстрационный прототип, связывающий экспериментальную модель с прикладным сценарием динамического изменения группы пользователей, а также формулируются направления дальнейшего развития технологии.

Важным ограничением работы является отсутствие доступа к внутренним данным и промышленным алгоритмам «Яндекс Музыки», а также к коммерческим и персональным данным пользователей. Поэтому промышленная часть сервиса описывается как концепт, а техническая часть проверяется на открытом обезличенном датасете YAMBDA. В этом датасете треки и пользователи анонимизированы, поэтому эксперимент демонстрирует качество модели через метрики рекомендаций, а не через интерпретируемые названия конкретных композиций.

% TODO(групповая ВКР): в конце введения описать распределение задач между участниками проекта.
